\newpage
\subsection{The finite linear gate across two consecutive layers}

\begin{result}{E.69}{Proved}{A finite linear gate across each odd--even pair of layers}
\label{coh:two-step-finite-gate}
Let an ordinary pair have support $\Sigma$ containing two
and satisfying
\[
 q\equiv-1\pmod3,\qquad v_3(q^2-1)=1
                       \quad(q\in\Sigma).
\]
The pairs of E.125 satisfy these hypotheses.
Fix an even integer $w\geq4$ and a global integral
pro-three class-$(w-2)$ section $s$ which is
boundary matched at every $q\in\Sigma$.
Put
\[
 b=b_{w-1},\qquad c=b_w,\qquad
 s_w=1+v_3(w),\qquad R_w=\mathbb Z/3^{s_w},
 \qquad C_w=R_w^c .
\]
Once a reference lift and cocycle bases are chosen,
there are defined matrices
$A_q\in\operatorname{Mat}_{c\times b}(R_w)$
and vectors $v_q\in C_w$, for $q\in\Sigma$, such
that a boundary-matched class-$w$ section above
$s$ exists if and only if the simultaneous system
\[
                         A_q\lambda=v_q
             \quad(q\in\Sigma\setminus\{2\})
 \tag{E.69.1}
\]
has a solution $\lambda\in R_w^b$.
For a fixed integral choice of the intervening
odd-layer parameters, matching at two fixes the
unique even-layer global twist.

Some set of at most $b+2$ support places, including
two, determines this entire matching condition.
Thus failure has a witness of that bounded size.
For $w=4$ there are exactly nine possible
odd-layer parameter classes to test, namely
$\lambda\in\mathbb F_3^2$, and at most four
support places suffice.

The statement does not assert consistency of
(E.69.1), an effective stopping rule for locating
the controlling places, or that higher digits
of the odd parameters are irrelevant to subsequent
pairs of layers.
\end{result}
\begin{proof}
Write $N_i=\Pi^{(3)}/\Gamma_{i+1}\Pi^{(3)}$.
The kernel of $N_w\to N_{w-2}$ is
\[
 M=\Gamma_{w-1}\Pi^{(3)}/\Gamma_{w+1}\Pi^{(3)}.
\]
It is abelian, since
$[\Gamma_{w-1},\Gamma_{w-1}]
\subseteq\Gamma_{2w-2}\subseteq\Gamma_{w+1}$.
Integral Hall coordinates identify it with a
finite free $\mathbb Z_3$-module.
The arithmetic action twisted by the section
$s$ makes it a continuous $G_{\mathbb Q}$-module.
To define it, lift $s(g)$ to the arithmetic
class-$w$ group and conjugate $M$.
Two lifts differ by an element of $M$, whose
inner action on $M$ is trivial; the action is
therefore independent of the choice and is
a genuine group action. It fits into
\[
 0\longrightarrow L_w=\mathbb Z_3(w)^c
   \longrightarrow M
   \longrightarrow L_{w-1}=\mathbb Z_3(w-1)^b
   \longrightarrow0 .
 \tag{E.69.2}
\]
These are continuous modules with the actual
twisted action. Only their two graded actions
have been replaced by their Tate descriptions.

There is at least one global class-$w$ lift
$t$ of $s$. E.66 gives a class-$(w-1)$
lift; since $w-1$ is odd,
$H^1(G_q,\mathbb Z_3(w-1))=0$ at each
support prime. Every such odd lift is
therefore boundary matched, and E.66 gives
its class-$w$ lift. This works for every
odd lift of $s$, not only a selected one.

Normalize the projection to have the same
literal representative of $s$. Two class-$w$
lifts differ by a continuous $M$-valued
one-cocycle. Conjugacy preserving that
projection is exactly conjugacy by $M$:
a nontrivial projected geometric stabilizer
would have a first nonzero weight coordinate
fixed by all $\chi_3^i$, which is impossible
in a torsion-free Tate module.
Thus the geometric-conjugacy classes form
the torsor $t+H^1(G_{\mathbb Q},M)$.
The local stabilizer argument at a support
prime uses the nonzero integer $q^i-1$
instead.

The long exact sequence of (E.69.2) and
$H^0(G_{\mathbb Q},L_{w-1})=0$ make
$H^1(G_{\mathbb Q},L_w)$ inject into
$H^1(G_{\mathbb Q},M)$.
Projection onto $H^1(G_{\mathbb Q},L_{w-1})$
is surjective: each class in the latter
group specifies an odd lift of $s$, and
the preceding argument supplies a global
class-$w$ lift above it. Subtract $t$ to
obtain an $M$-cocycle with that projection.
This uses the proved existence theorem;
it does not assert that the absolute group
$H^2(G_{\mathbb Q},L_w)$ is zero.
E.67 consequently gives the exact sequence
\[
 0\longrightarrow C_w
 \xrightarrow{\ \iota\ } H^1(G_{\mathbb Q},M)
 \longrightarrow\mathbb Z_3^b
 \longrightarrow0 .
 \tag{E.69.3}
\]

Choose an $M$-valued cocycle lifting each
basis vector of the last free module.
Their linear combinations give a splitting
$j:\mathbb Z_3^b\to H^1(G_{\mathbb Q},M)$.
It is continuous in the continuous-cochain
topology: each chosen cocycle has compact
image, and multiplying its values by
convergent three-adic scalars gives uniform
convergence. A finite sum is therefore a
continuous family of continuous cocycles.
The finite central factor is represented
by the cocycles $\eta_w$ of E.67.
Every global lift class is uniquely of
the form
\[
            t+\iota(\tau)+j(\lambda),
       \qquad \tau\in C_w,\quad
              \lambda\in\mathbb Z_3^b .
 \tag{E.69.4}
\]
No rational splitting or extra ramification
in arbitrary finite cocycles is used.

At $q\in\Sigma$, the odd weight $w-1$
has both zero invariants and zero first
cohomology. The local long exact sequence
therefore gives
\[
 H^1(G_q,L_w)\ \xrightarrow{\ \sim\ }\
                         H^1(G_q,M).
 \tag{E.69.5}
\]
Also $v_3(q^w-1)=1+v_3(w)=s_w$.
By E.67, restriction
$C_w=H^1(G_{\mathbb Q},L_w)\to
H^1(G_q,L_w)$ is an isomorphism.
In each coordinate its generator has
Frobenius value $(q^w-1)/3^{s_w}$,
with the corresponding unit factor
retained. Combining these maps gives
$k_q:C_w\xrightarrow{\sim}H^1(G_q,M)$.

Align the prescribed actual local point
below class $w$ with $s|_{G_q}$.
The projected conjugator is unique by
the stabilizer argument; different lifts
of it differ by $M$ and hence change
the local lift only by a coboundary.
The local point therefore has a
well-defined discrepancy from $t$ in
$H^1(G_q,M)$. Let $d_q\in C_w$ be
its inverse image under $k_q$, and
define the $\mathbb Z_3$-linear map
\[
 B_q:\mathbb Z_3^b\longrightarrow C_w,
 \qquad B_q=k_q^{-1}\circ
                    \operatorname{res}_q\circ j .
\]
The class (E.69.4) matches that point
precisely when
\[
                         \tau+B_q\lambda=d_q .
 \tag{E.69.6}
\]
Every $B_q$ factors through $R_w^b$,
since its target is killed by $3^{s_w}$.
At the support place two, (E.69.6)
determines $\tau=d_2-B_2\lambda$.
Substitution elsewhere gives (E.69.1)
with $A_q=B_q-B_2$ and $v_q=d_q-d_2$.
Every solution modulo $3^{s_w}$ may
be lifted to integral odd parameters;
with this central twist it yields an
actual matched class-$w$ section.

For the bounded-place assertion,
treat every scalar equation in
(E.69.1) as an augmented row
$(a_1,\ldots,a_b,v)\in R_w^{b+1}$.
Let $V$ be their generated submodule.
It has at most $b+1$ generators:
Smith normal form over the finite
principal ideal ring $R_w$ gives
this bound for a submodule of its
free module of rank $b+1$.
The generators can be chosen from
the original rows. Choose a basis
of $V/3V$ from their images, then
use Nakayama's lemma.
Every other augmented row is an
$R_w$-linear combination of the
chosen rows. Satisfying their
equations consequently implies
all the equations, including the
right sides. They involve at most
$b+1$ nonpinning support places;
adding two gives at most $b+2$.

In weight four the free Lie algebra
on two generators has weight-three
basis $[x,[x,y]]$, $[y,[x,y]]$.
Thus $b_3=2$ and $s_4=1$:
the parameter space is $\mathbb F_3^2$
and the place bound is four.
The finite controlling subset exists
because its entire augmented-row
module is finite. This gives no
effective test that a presently
enumerated subset already generates
all remaining rows, and gives no
assertion that the resulting equations
are consistent.
\end{proof}
